\subsection{基数之间的序关系}

关系“X与Y的某个子集等势”等价于“存在从X到Y的单射”；它也等价于关系“Card(X)与Card(Y)的某个子集等势”（第二章，§3，第8号，定理1）。

定理1. 关系 \(R\{\mathfrak{x}, \mathfrak{y}\}\) : `` \(\mathfrak{x}\) 与 \(\mathfrak{y}\) 是基数且 \(\mathfrak{x}\) 与 \(\mathfrak{y}\) 的某个子集等势 '' 是一个良序关系（§2，第1号）。

由于 \(\mathbf{R}\{\mathfrak{x},\mathfrak{x}\}\) 对每个基数 \(\mathfrak{x}\) 都成立，那么需要证明的是，对每个由基数组成的集合 \(\mathbf{E}\) ，关系 `` \(\mathfrak{x} \in \mathbf{E}\) 并且 \(\mathfrak{y} \in \mathbf{E}\) 并且 \(\mathbf{R}\{\mathfrak{x},\mathfrak{y}\}\) '' 是 \(\mathbf{E}\) 上的一个良序关系。考虑集合 \(\mathbf{A} = \bigcup_{\mathfrak{r}\in \mathbf{E}}\mathfrak{x}\) .

每个基数 \(\mathfrak{r} \in \mathbf{E}\) 都是 A 的子集。由 §2 定理 1，A 上存在一个良序关系，我们将其记为 \(\mathfrak{r} \leqslant y\) ，并且 A 的每个子集都与 A 的一个段等势（§2，第5号，定理 3，推论 \(x\) ）。对每个基数 \(\mathfrak{r} \in \mathbf{E}\) ，考虑 A 中与 \(\mathfrak{r}\) 等势的段的集合；这个段的集合非空，因此有一个最小元素（§2，第1号，命题 2）；将此最小元素记为 \(\varphi(\mathfrak{r})\) 。关系

`` \(x \in E\) 且 \(y \in E\) 且 \(x\) 与 \(y\) 的某个子集等势 ''

则等价于

\[
" \mathfrak {x} \in \mathrm{E} \text {   且   } \mathfrak {y} \in \mathrm{E} \text {   且   } \varphi (\mathfrak {x}) \subset \varphi (\mathfrak {y})".
\]

因为显然第二个关系蕴含第一个。反之，如果 \(\mathfrak{r}\) 与 \(\varphi(\mathfrak{y})\) 的某个子集等势，我们不能有 \(\varphi(\mathfrak{y}) \subset \varphi(\mathfrak{x})\) 且 \(\varphi(\mathfrak{y}) \neq \varphi(\mathfrak{x})\) ；否则会存在 \(\varphi(\mathfrak{y})\) 的一段与 \(\mathcal{X}\) 等势 ( \(\S\) 2，第5号，定理3，推论3），这与 \(\varphi(\mathfrak{x})\) 的定义相矛盾。由于 A 的段的集合按包含关系是良序的（ \(\S\) 2，第1号，命题2），定理得证。

我们将关系 \(\mathbf{R}\{\mathfrak{x},\mathfrak{y}\}\) 记为 \(\mathfrak{x}\leqslant \mathfrak{y}\) 。一个集合 \(\mathbf{X}\) 与某个集合 \(\mathbf{Y}\) 的子集等势，当且仅当 \(\operatorname{Card}(\mathbf{X})\leqslant \operatorname{Card}(\mathbf{Y})\) .

显然我们有 \(0 \leqslant \mathfrak{r}\) （对于每一个基数 \(\mathfrak{r}\) ），以及 \(1 \leqslant \mathfrak{r}\) （对于每一个基数 \(\mathfrak{r} \neq 0\) ）.

推论1. 给定任意两个集合，其中一个集合与另一个集合的某个子集等势。

推论2. 两个集合，若每一个都与另一个的子集等势，则它们等势。

注记. 给定任意集合A，存在一个集合，其元素为基数 \(\operatorname{Card}(X)\) （对于A的所有子集X），即形如 \(\operatorname{Card}(X)\) （对于 \(X \in \mathfrak{P}(A)\) ）的客体所组成的集合（第二章，§1，第6号）。对于每个基数a，关系“r是一个基数且 \(r \leqslant a\) ”因而在 r 中是收集性的（第二章，§1，第4号），因为它等价于关系“ \(\mathfrak{f}\) 具有 Card(X) 的形式（其中 \(X \subset \mathfrak{a}\) ）”；所有满足此关系的 \(\mathfrak{f}\) 的集合称为基数 \(\leqslant \mathfrak{a}\) 的集合.

命题2. 对于每个由基数组成的族 \((\mathfrak{a}_i)_{i\in I}\) ，存在唯一的基数 \(\mathfrak{b}\) ，使得 \(\mathfrak{a}_i\leqslant \mathfrak{b}\) （对于所有 \(i\in I\) ），并且使得每个基数 \(\mathfrak{c}\) （满足 \(\mathfrak{a}_i\leqslant \mathfrak{c}\) ，对于所有 \(i\in I\) ）都 \(\geqslant \mathfrak{b}\) .

存在一个包含所有集合 \(a_{i}\) 的集合 E（例如，这些集合的和（第二章，§4，第8号）），由此 \(a_{i} \leqslant a = Card(E)\) （对于所有 \(i \in I\) ）。基数 \(\leqslant a\) 的集合 F 是良序的，并且包含所有的 \(a_{i}\) ，因此族 \((a_{i})_{i \in I}\) 在 F 中有最小上界 b。设 c 为一个基数，它 \(\geqslant a_{i}\) （对于所有 \(i \in I\) ）；若 \(c < b \leqslant a\) ，那么 \(c \in F\) ，而不等式 \(a_{i} \leqslant c\) 与族 \((a_{i})\) 在有序集 F 中最小上界的定义相矛盾；因此结论成立。

借用语言习惯，基数 \(\mathfrak{b}\) 称为族 \((\mathfrak{a}_i)_{i\in I}\) 的最小上界，并且用 \(\sup_{i\in I}a_i\) 表示.

命题3. 设X和Y为集合。若存在从X到Y的满射f，则Card(Y) ≤ Card(X)。

因为存在一个截面 \(s\) （与 \(f\) 相关联）(第二章，§3，第8号，命题8)，且 \(s\) 是Y到X的一个单射。

